\chapter{Topologi Ruang Metrik}\label{ch:b2:metric}

Topologi garis real pada jilid Tahun ke-1 menyeluruh, nyaris
tanpa mengubah satu kata pun, ke sembarang himpunan yang dilengkapi
jarak. Untungnya besar sekali: barisan fungsi, matriks, kurva --- semuanya
menjadi titik \omterm{def:b2:metric:def}{ruang metrik}, dan ketiga pilar yang dibuktikan di sini
--- kelengkapan beserta teorema titik tetap Banach, kekompakan,
keterhubungan --- berlaku bagi semuanya secara seragam. Bab ini adalah
tulang punggung seluruh paruh analisis buku ini.

\section{Ruang metrik}

\begin{definition}\label{def:b2:metric:def}
Sebuah \emph{ruang metrik}\index{ruang metrik} adalah himpunan $X$ beserta
pemetaan $d \colon X \times X \to \R_+$ sehingga, untuk setiap $x, y, z$:
\[
d(x,y) = 0 \iff x = y,
\qquad
d(x,y) = d(y,x),
\qquad
d(x,z) \leq d(x,y) + d(y,z).
\]
Bolanya: $B(a, r) = \{x : d(a,x) < r\}$ (\omterm{def:b2:metric:topology}{terbuka}), $\overline B(a,r) =
\{x : d(a,x) \leq r\}$ (tertutup). Sebuah himpunan bagian $A \subseteq X$
menjadi ruang metrik dengan jarak imbasannya.
\end{definition}

\begin{example}\label{ex:b2:metric:examples}
$\R$ dengan $\abs{x - y}$; $\R^n$ dengan salah satu dari
\[
d_1(x,y) = \sum_i \abs{x_i - y_i},
\quad
d_2(x,y) = \Bigl(\sum_i (x_i - y_i)^2\Bigr)^{1/2},
\quad
d_\infty(x,y) = \max_i \abs{x_i - y_i};
\]
himpunan $C(\intcc{a}{b})$ berisi fungsi \omterm{def:b2:metric:continuity}{kontinu} dengan
\emph{jarak supremum} $d_\infty(f, g) = \sup_{\intcc{a}{b}} \abs{f -
g}$ (yang hingga, sebab $f - g$ terbatas); sembarang himpunan dengan
jarak \emph{diskret} ($d(x,y) = 1$ untuk $x \neq y$). Jarak yang berasal
dari norma menjadi bahasan \cref{ch:b2:nvs}.
\end{example}

\begin{definition}[Topologi sebuah ruang metrik]\label{def:b2:metric:topology}
$U \subseteq X$ disebut \emph{terbuka}\index{himpunan terbuka!pada ruang metrik}
apabila setiap titik $U$ menjadi pusat sebuah bola yang termuat di $U$;
sedangkan $F$ disebut \emph{tertutup} apabila komplemennya terbuka.
Persekitaran, interior, penutup, kepadatan, dan batas ditetapkan persis
seperti pada garis real (jilid Tahun ke-1), dengan bola menggantikan
interval, dan pernyataan yang dibuktikan di sana --- gabungan dan irisan
himpunan terbuka, pencirian interior dan penutup, penutup sebagai
himpunan tertutup terkecil yang memuatnya --- berpindah dengan bukti yang
sama. Bola terbuka memang terbuka, dan bola tertutup memang tertutup
(lewat ketaksamaan segitiga).
\end{definition}

\begin{example}[Interior, penutup, batas pada satu himpunan]\label{ex:b2:metric:intclosure}
Di $\R$, misalkan $A = \intoc{0}{1} \cup \{2\}$. Interiornya:
$\intoo{0}{1}$ --- di sekitar sembarang $x \in \intoo01$ ada bola kecil
yang tetap di $A$; sedangkan di sekitar $1$, setiap bola $\intoo{1-r}{1+r}$
bocor keluar dari $A$ di sebelah kanan, jadi $1$ bukan titik interior;
dan titik terpencil $2$ pun bukan titik interior. Penutupnya:
$\intcc{0}{1} \cup \{2\}$ (titik $0$ adalah limit $A$, dan tak ada
lagi yang bertambah). Batasnya (penutup dikurangi interior):
$\{0, 1, 2\}$. Perhatikan ketidaksetangkupan yang patut diingat: sebuah
ujung dapat menjadi anggota himpunan tanpa menjadi titik interior (yaitu
$1$), dapat melekat tanpa menjadi anggota (yaitu $0$), dan titik
terpencil menjadi batas bagi dirinya sendiri (yaitu $2$). Pembukuan yang
sama berjalan kata demi kata pada sembarang \omterm{def:b2:metric:def}{ruang metrik}, dengan bola
menggantikan interval.
\end{example}

\begin{definition}[Limit, kekontinuan]\label{def:b2:metric:continuity}
Kita tulis $x_n \to x$ di $X$ apabila $d(x_n, x) \to 0$. Pemetaan $f
\colon X \to Y$ antara \omterm{def:b2:metric:def}{ruang metrik} disebut
\emph{kontinu}\index{kontinu} di $a$ apabila
\[
\forall \varepsilon > 0,\ \exists\delta > 0,\quad
d_X(x, a) \leq \delta \implies d_Y\bigl(f(x), f(a)\bigr) \leq
\varepsilon ;
\]
setara dengan itu (buktinya sama seperti pada $\R$), $f(x_n) \to f(a)$
untuk setiap barisan $x_n \to a$. Adapun $f$ disebut
\emph{Lipschitz}\index{Lipschitz} dengan konstanta $k$ apabila selalu
berlaku $d_Y(f(x), f(y)) \leq k\, d_X(x, y)$ --- dan ketika itu ia
kontinu seragam, jadi kontinu.
\end{definition}

\begin{theorem}[Pencirian global kekontinuan]\label{thm:b2:metric:globalcontinuity}
$f \colon X \to Y$ \omterm{def:b2:metric:continuity}{kontinu} (di setiap titik) bila dan hanya bila
prapeta setiap himpunan \omterm{def:b2:metric:topology}{terbuka} adalah \omterm{def:b2:metric:topology}{terbuka} --- bila dan hanya bila
prapeta setiap himpunan tertutup adalah tertutup.
\end{theorem}

\begin{proof}
($\Rightarrow$) Misalkan $V \subseteq Y$ \omterm{def:b2:metric:topology}{terbuka} dan $a \in f^{-1}(V)$:
ada bola $B(f(a), \varepsilon) \subseteq V$; kekontinuan di $a$
menyediakan $\delta$ dengan $f(B(a, \delta)) \subseteq B(f(a),
\varepsilon)$, sehingga $B(a, \delta) \subseteq f^{-1}(V)$.

($\Leftarrow$) Diberikan $a$ dan $\varepsilon$: himpunan
$f^{-1}\bigl(B(f(a), \varepsilon)\bigr)$ \omterm{def:b2:metric:topology}{terbuka} dan memuat $a$, jadi ia
memuat sebuah bola $B(a, \delta)$: dan itulah definisi kekontinuan di
$a$. Untuk himpunan tertutup: ambil komplemennya
(\cref{prop:b2:structures:images}).
\end{proof}

\section{Ruang lengkap}

\begin{definition}\label{def:b2:metric:complete}
Sebuah barisan $(x_n)$ disebut \emph{Cauchy} apabila $\sup_{p, q \geq N}
d(x_p, x_q) \to 0$ ketika $N \to \infty$. Sebuah \omterm{def:b2:metric:def}{ruang metrik} disebut
\emph{lengkap}\index{ruang metrik lengkap} apabila setiap barisan Cauchy
konvergen. Konvergen $\Rightarrow$ Cauchy selalu berlaku; himpunan bagian
tertutup dari ruang lengkap bersifat lengkap, dan himpunan bagian lengkap
dari ruang apa pun bersifat tertutup (buktinya sama seperti pada $\R$:
jilid Tahun ke-1).
\end{definition}

\begin{example}[Cauchy tanpa limit]\label{ex:b2:metric:cauchynolimit}
Di $X = \Q$ dengan jarak biasa, pemotongan desimal
$\sqrt2$,
\[
x_0 = 1,\quad x_1 = 1.4,\quad x_2 = 1.41,\quad x_3 = 1.414,
\quad\dots
\]
memenuhi $\abs{x_p - x_q} \leq 10^{-\min(p,q)}$: jadi Cauchy di $\Q$.
Limitnya di $\Q$ akan sekaligus menjadi limitnya di $\R$, yakni
$\sqrt2 \notin \Q$: jadi tak ada limitnya di $X$. Ketaklengkapan adalah
hadirnya ``limit hantu'' semacam itu; kelengkapan $\R$
dirancang pada jilid Tahun ke-1 justru untuk memberi rumah bagi setiap
barisan Cauchy.
\end{example}

\begin{theorem}\label{thm:b2:metric:rncomplete}
Ruang $\R^n$ (dengan salah satu dari ketiga jarak pada
\cref{ex:b2:metric:examples}) dan $\bigl(C(\intcc{a}{b}),
d_\infty\bigr)$ bersifat \omterm{def:b2:metric:complete}{lengkap}.
\end{theorem}

\begin{proof}
Untuk $\R^n$: barisan Cauchy bersifat Cauchy pada tiap koordinatnya
(sebab $\abs{x_i - y_i} \leq d(x,y)$ untuk ketiga jaraknya), jadi tiap
koordinatnya konvergen (kelengkapan $\R$, jilid Tahun ke-1), dan
kekonvergenan koordinat demi koordinat mengakibatkan kekonvergenan untuk
$d_\infty$ (sebab koordinatnya berhingga banyak), jadi untuk ketiganya
(ketiga jarak itu saling mendominasi dalam faktor konstan: $d_\infty
\leq d_2 \leq d_1 \leq n\,d_\infty$).

Untuk $C(\intcc{a}{b})$: misalkan $(f_n)$ bersifat Cauchy terhadap
$d_\infty$. Untuk tiap $x$, barisan $(f_n(x))$ Cauchy di $\R$ (sebab
$\abs{f_p(x) - f_q(x)} \leq d_\infty(f_p, f_q)$), jadi ia konvergen ke
suatu $f(x)$. Dengan melewatkan ke limit pada $\abs{f_p(x) -
f_q(x)} \leq \varepsilon$ (yang berlaku untuk $p, q \geq N_\varepsilon$
dan setiap $x$) ketika $q \to \infty$: diperoleh $\abs{f_p(x) - f(x)} \leq \varepsilon$ untuk
setiap $x$, yakni $d_\infty(f_p, f) \leq \varepsilon$: itulah kekonvergenan
seragam. Limitnya \omterm{def:b2:metric:continuity}{kontinu}: diberikan $\varepsilon$, pilih $p$ dengan
$\sup\abs{f_p - f} \leq \varepsilon$, lalu pakai kekontinuan $f_p$ di $a$
beserta pemecahan tiga suku
\[
\abs{f(x) - f(a)} \leq \abs{f(x) - f_p(x)} + \abs{f_p(x) - f_p(a)}
+ \abs{f_p(a) - f(a)} \leq 3\varepsilon
\]
untuk $x$ yang dekat ke $a$. (``Hujah $3\varepsilon$'' ini kembali
sebagai teorema limit seragam pada \cref{ch:b2:funcseq}.)
\end{proof}

\begin{example}[Himpunan terbuka dan tertutup dikenali lewat kekontinuan]\label{ex:b2:metric:recognize}
Pencirian global
(\cref{thm:b2:metric:globalcontinuity}) adalah perkakas sehari-hari untuk
pembukuan topologis. Di $\R^2$: himpunan $\{(x, y) : x^2 + y^2
< 1,\ y > x^3\}$ \omterm{def:b2:metric:topology}{terbuka} --- sebab ia $g^{-1}(\intoo{-\infty}{1})
\cap h^{-1}(\intoo{0}{+\infty})$ untuk $g(x,y) = x^2 + y^2$ dan
$h(x, y) = y - x^3$ yang \omterm{def:b2:metric:continuity}{kontinu}, yakni irisan dua prapeta
\omterm{def:b2:metric:topology}{terbuka}. Di $\bigl(C(\intcc01), d_\infty\bigr)$: himpunan semua
fungsi dengan $f(0) = f(1)$ dan $\int_0^1 f = 0$ bersifat tertutup ---
sebab ia prapeta $\{(0,0)\}$ di bawah pemetaan \omterm{def:b2:metric:continuity}{kontinu} $f \mapsto
\bigl(f(0) - f(1),\ \int_0^1 f\bigr)$ ke $\R^2$ (tiap
koordinatnya \omterm{def:b2:metric:continuity}{Lipschitz} dengan konstanta $1$, seperti pada
\cref{exo:b2:metric:3}). Metodenya tak pernah menggambar apa pun:
tunjukkan sebuah pemetaan \omterm{def:b2:metric:continuity}{kontinu}, bacalah himpunannya sebagai prapeta,
lalu kutip teoremanya.
\end{example}

\begin{example}[Himpunan tertutup yang ditetapkan oleh tak hingga banyak syarat]\label{ex:b2:metric:lipclosed}
Di $\bigl(C(\intcc01), d_\infty\bigr)$, himpunan
\[
L = \{f : \abs{f(x) - f(y)} \leq \abs{x - y}
\ \text{untuk setiap } x, y\}
\]
berisi fungsi \omterm{def:b2:metric:continuity}{Lipschitz} berkonstanta $1$ bersifat tertutup, walaupun ia
dipahat oleh tak \omterm{def:b2:structures:countable}{terbilang} banyak syarat: untuk tiap pasangan tetap $(x,
y)$, pemetaan $f \mapsto \abs{f(x) - f(y)} - \abs{x - y}$ bersifat
\omterm{def:b2:metric:continuity}{kontinu} (sebab evaluasi bersifat \omterm{def:b2:metric:continuity}{Lipschitz} berkonstanta $1$), jadi tiap
syarat tunggalnya menetapkan himpunan tertutup, dan $L$ adalah
\emph{irisan} keluarga itu --- sedangkan irisan sembarang atas himpunan
tertutup tetap tertutup. Pola yang sama menjamin ketertutupan bagi fungsi
monoton, fungsi cembung, dan fungsi yang terbatas oleh $g$ tertentu:
limit seragam mewarisi setiap sifat yang dapat diungkapkan sebagai
keluarga kendala titik demi titik yang tertutup. Adapun apa yang
\emph{tidak} otomatis diwarisi limit seragam --- keterdiferensialan,
misalnya --- justru itulah yang harus diperjuangkan \cref{ch:b2:funcseq}.
\end{example}

\begin{theorem}[Teorema titik tetap Banach]\label{thm:b2:metric:banach}
Misalkan $X$ \omterm{def:b2:metric:complete}{ruang metrik lengkap} yang tak kosong dan $f \colon X \to X$
sebuah \emph{kontraksi}: \omterm{def:b2:metric:continuity}{Lipschitz} dengan konstanta $k < 1$. Maka $f$
punya tepat satu titik tetap $\ell$, dan setiap orbit $x_{n+1} = f(x_n)$
konvergen ke $\ell$, dengan\index{teorema titik tetap Banach}
\[
d(x_n, \ell) \leq \frac{k^n}{1 - k}\, d(x_1, x_0) .
\]
\end{theorem}

\begin{proof}
Ketunggalan: jarak dua titik tetap bernilai $\leq k$ kali jarak itu
sendiri. Keberadaan: secara induktif $d(x_{n+1}, x_n) \leq k^n d(x_1,
x_0)$, sehingga untuk $q > p$,
\[
d(x_q, x_p) \leq \sum_{j=p}^{q-1} d(x_{j+1}, x_j)
\leq d(x_1, x_0) \sum_{j \geq p} k^j
= \frac{k^p}{1-k}\, d(x_1, x_0) \xrightarrow[p\to\infty]{} 0 :
\]
jadi Cauchy, sehingga konvergen ke suatu $\ell$; kekontinuan $f$
melewatkan $x_{n+1} = f(x_n)$ ke limitnya: $\ell = f(\ell)$. Batas
galatnya adalah taksiran di atas dengan $q \to \infty$.
\end{proof}

\begin{example}[Sebuah persamaan integral]\label{ex:b2:metric:integralequation}
Pada $X = C(\intcc{0}{1})$ (yang \omterm{def:b2:metric:complete}{lengkap},
\cref{thm:b2:metric:rncomplete}), tinjaulah $T(f)(x) = 1 + \frac12
\int_0^x f(t)\,\dd t$. Untuk $f, g \in X$:
\[
\abs{T(f)(x) - T(g)(x)} \leq \frac12 \int_0^x \abs{f - g} \leq
\frac12\, d_\infty(f, g),
\]
jadi $T$ adalah kontraksi berkonstanta $\frac12$: ia punya tepat satu
titik tetap yang \omterm{def:b2:metric:continuity}{kontinu} --- yaitu penyelesaian $f' = \frac f2$, $f(0) =
1$, yakni $\eu^{x/2}$. Skema ini, setelah diindustrikan, menjadi
teorema Cauchy--Lipschitz pada \cref{ch:b2:diffeq}.
\end{example}

\begin{example}[Sebuah titik tetap numerik: $x = \cos x$]\label{ex:b2:metric:cosfix}
Pada $X = \intcc{0}{1}$ yang \omterm{def:b2:metric:complete}{lengkap}, pemetaan $f = \cos$ mengirim $X$
ke dalam $\intcc{\cos 1}{1} \subseteq X$ dan merupakan kontraksi: menurut
ketaksamaan nilai rata-rata,
\[
\abs{\cos x - \cos y} \leq \bigl(\sup_{\intcc01}\abs{\sin}\bigr)
\abs{x - y} = (\sin 1)\abs{x - y},
\qquad \sin 1 \approx 0.841 < 1 .
\]
Menurut Banach: ada tepat satu penyelesaian $x = \cos x$ di $\intcc01$
(jadi juga di $\R$: sebab setiap titik tetap real terletak di
$\intcc{-1}{1}$, lalu di $\intcc{\cos 1}{1}$ setelah satu penerapan), dan
iterasi $x_{n+1} = \cos x_n$ konvergen ke sana dari awal mana pun:
$x_\infty \approx 0.739085$, bilangan masyhur yang diperoleh dengan
menggenjot tombol kosinus sebuah kalkulator. Batas galatnya meramalkan
peluruhan $(\sin1)^n/(1 - \sin1)$ --- kira-kira satu angka tiap $13$
tekanan; adapun batas aposteriori pada soal akhir pekan bab
ini (pertanyaan 14) mengesahkan tiap langkahnya secara langsung.
\end{example}

\section{Kekompakan}

\begin{definition}\label{def:b2:metric:compact}
Sebuah \omterm{def:b2:metric:def}{ruang metrik} $X$ disebut \emph{kompak}\index{ruang metrik kompak}
apabila setiap barisan di $X$ punya barisan bagian yang konvergen
\emph{di dalam} $X$ (sifat Bolzano--Weierstrass). Sebuah himpunan bagian
disebut kompak apabila ia kompak dengan jarak imbasannya.
\end{definition}

\begin{theorem}[Sifat pertama]\label{thm:b2:metric:compactprops}
\begin{enumerate}
  \item Himpunan bagian yang \omterm{def:b2:metric:compact}{kompak} bersifat tertutup dan terbatas;
        himpunan bagian tertutup dari ruang \omterm{def:b2:metric:compact}{kompak} bersifat \omterm{def:b2:metric:compact}{kompak}.
  \item Di $\R^n$ konversnya berlaku: \omterm{def:b2:metric:compact}{kompak} $\iff$ tertutup dan
        terbatas.
  \item Peta \omterm{def:b2:metric:continuity}{kontinu} sebuah ruang \omterm{def:b2:metric:compact}{kompak} bersifat \omterm{def:b2:metric:compact}{kompak}; dan
        fungsi real \omterm{def:b2:metric:continuity}{kontinu} pada ruang \omterm{def:b2:metric:compact}{kompak} tak kosong bersifat
        terbatas serta mencapai batasnya.
  \item (Heine) Pemetaan \omterm{def:b2:metric:continuity}{kontinu} pada ruang \omterm{def:b2:metric:compact}{kompak} bersifat kontinu
        seragam.
  \item Hasil kali: jika $X, Y$ \omterm{def:b2:metric:compact}{kompak}, maka $X \times Y$ \omterm{def:b2:metric:compact}{kompak} (dengan
        $d\bigl((x,y),(x',y')\bigr) = d(x,x') + d(y,y')$).
\end{enumerate}
\end{theorem}

\begin{proof}
(1) Hujahnya sama seperti pada garis real (jilid Tahun ke-1): barisan yang
lolos ke tak hingga atau yang konvergen ke luar tidak punya barisan bagian
yang konvergen di dalam; untuk pernyataan kedua, sarikan barisan bagiannya
di dalam ruang \omterm{def:b2:metric:compact}{kompak} sekitarnya lalu pakai ketertutupannya.

(2) Barisan terbatas di $\R^n$ punya barisan bagian yang konvergen
komponen demi komponen: sarikan pada koordinat pertama
(Bolzano--Weierstrass pada $\R$), lalu dari barisan bagian itu sarikan
pada koordinat kedua, dan seterusnya ($n$ penyarian berturut-turut);
ketertutupannya menjaga limitnya tetap di dalam.

(3) Diberikan $(f(x_n))$, sarikan $x_{\varphi(n)} \to x \in X$;
kekontinuan memberi $f(x_{\varphi(n)}) \to f(x) \in f(X)$. Untuk kasus
realnya: kekompakan $f(X) \subseteq \R$ membuatnya tertutup dan terbatas,
dan $\sup f(X) \in f(X)$ (sebab supremum sebuah himpunan melekat padanya,
dan $f(X)$ tertutup).

(4) Bukti Tahun ke-1 berpindah kata demi kata; inilah bunyinya dalam
busana metrik. Andaikan $f \colon X \to Y$ \omterm{def:b2:metric:continuity}{kontinu} pada $X$ yang \omterm{def:b2:metric:compact}{kompak}
tetapi tidak kontinu seragam: maka ada $\varepsilon > 0$ yang, untuk
setiap $n$, menyediakan titik dengan
\[
d_X(x_n, y_n) \leq \frac{1}{n+1}
\qquad\text{dan}\qquad
d_Y\bigl(f(x_n), f(y_n)\bigr) > \varepsilon .
\]
Sarikan $x_{\varphi(n)} \to a \in X$; maka $y_{\varphi(n)} \to
a$ juga (sebab jarak keduanya menuju $0$). Kekontinuan di $a$
mengirim kedua barisan petanya ke $f(a)$, sehingga
$d_Y\bigl(f(x_{\varphi(n)}), f(y_{\varphi(n)})\bigr) \to 0$ ---
dan itu bertentangan dengan jurang seragam $> \varepsilon$. Kekompakan
menyediakan tepat satu hal: titik kumpul $a$, tempat kekontinuan biasa
diterapkan.

(5) Sarikan pada koordinat $X$, lalu sarikan lagi pada koordinat
$Y$.
\end{proof}

\begin{example}[Teorema Heine, dengan dan tanpa kekompakan]\label{ex:b2:metric:heinework}
Pada $\intcc{0}{1}$, fungsi $x \mapsto x^2$ bersifat kontinu
seragam --- Heine mengatakannya tanpa perhitungan, tetapi taksiran
langsungnya tetap mendidik:
\[
\abs{x^2 - y^2} = \abs{x + y}\,\abs{x - y} \leq 2\abs{x - y},
\]
jadi $\delta = \varepsilon/2$ berhasil \emph{untuk semua titik
sekaligus}. Pada $\R$ fungsi yang sama tidak kontinu seragam:
dengan $x_n = n$ dan $y_n = n + \frac1n$, jurangnya $\abs{x_n - y_n}
= \frac1n \to 0$ padahal $\abs{x_n^2 - y_n^2} = 2 + \frac1{n^2}
\geq 2$: jadi tak ada satu pun $\delta$ yang melayani $\varepsilon = 1$.
Mekanismenya kasatmata: konstanta \omterm{def:b2:metric:continuity}{Lipschitz} lokalnya $\abs{x + y}$
terbatas pada himpunan \omterm{def:b2:metric:compact}{kompak} dan tak terbatas pada $\R$ --- teorema
Heine adalah pernyataan bahwa kekompakan membatasi konstanta lokal
semacam itu secara seragam.
\end{example}

\begin{method}[Membuktikan sebuah himpunan kompak]\label{met:b2:metric:compactness}
Ada tiga jalan, menurut seringnya dipakai. (1) \emph{Pengenalan dari
ruang sekitarnya:} di $\R^n$ (atau di ruang bernorma berdimensi hingga
mana pun, \cref{ch:b2:nvs}), periksalah ketertutupannya --- lazimnya
sebagai prapeta, \cref{ex:b2:metric:recognize} --- dan keterbatasannya. (2)
\emph{Pewarisan:} himpunan bagian tertutup dari ruang \omterm{def:b2:metric:compact}{kompak} yang sudah
dikenal bersifat \omterm{def:b2:metric:compact}{kompak}; gabungan hingga atau hasil kali himpunan \omterm{def:b2:metric:compact}{kompak}
bersifat \omterm{def:b2:metric:compact}{kompak}; peta \omterm{def:b2:metric:continuity}{kontinu} himpunan \omterm{def:b2:metric:compact}{kompak} bersifat \omterm{def:b2:metric:compact}{kompak}. (3)
\emph{Tangan kosong:} sarikan barisan bagian yang konvergen dari barisan
sembarang --- biasanya lewat penyarian berturut-turut koordinat demi
koordinat. Untuk membuktikan ketidakkompakan, satu saksi sudah cukup:
sebuah barisan tanpa barisan bagian yang konvergen, paling sering berupa
titik yang jarak antaranya $\geq \varepsilon$.
\end{method}

\begin{example}[Jarak antara dua himpunan: kekompakan membayar sewanya]\label{ex:b2:metric:setdistance}
Misalkan $K$ \omterm{def:b2:metric:compact}{kompak}, $F$ tertutup, dan $K \cap F = \emptyset$ di dalam
sebuah \omterm{def:b2:metric:def}{ruang metrik}. Maka
\[
d(K, F) = \inf\,\{d(x, y) : x \in K,\ y \in F\} > 0 :
\]
sebab fungsi $x \mapsto d(x, F)$ bersifat \omterm{def:b2:metric:continuity}{kontinu}
(\cref{exo:b2:metric:11}) dan positif pada $K$ (sebab $d(x, F) = 0$
akan menaruh $x \in \overline F = F$), jadi ia mencapai minimum yang
positif pada himpunan \omterm{def:b2:metric:compact}{kompak} $K$
(\cref{thm:b2:metric:compactprops}~(3)). Kekompakan bukan
hiasan: untuk dua himpunan \emph{tertutup} infimumnya bisa lenyap
tanpa tercapai --- di $\R^2$, hiperbola $F_1 = \{xy
= 1\}$ dan sumbu $F_2 = \{y = 0\}$ merupakan dua himpunan tertutup yang
saling lepas dengan $d(F_1, F_2) = 0$ (sebab titik $(n, \frac1n)$
mendekati sumbunya). Justru pelarian ke tak hingga itulah yang dilarang kekompakan.
\end{example}

\begin{theorem}[Borel--Lebesgue]\label{thm:b2:metric:borellebesgue}
Sebuah \omterm{def:b2:metric:def}{ruang metrik} $X$ bersifat \omterm{def:b2:metric:compact}{kompak} bila dan hanya bila setiap
selimut $X$ oleh himpunan \omterm{def:b2:metric:topology}{terbuka} punya subselimut yang
\emph{hingga}.\index{teorema Borel--Lebesgue}
\end{theorem}

\begin{proof}
($\Leftarrow$) Andaikan $(x_n)$ tidak punya barisan bagian yang
konvergen. Kita klaim setiap $x \in X$ punya bola $B(x, r_x)$ yang memuat
$x_n$ hanya untuk berhingga banyak indeks $n$: sebab jika tidak, setiap
bola $B(x, \frac1{k+1})$ akan memuat tak hingga banyak sukunya, dan
memilih indeks
\[
\varphi(0) < \varphi(1) < \varphi(2) < \cdots
\quad\text{dengan}\quad
x_{\varphi(k)} \in B\Bigl(x, \frac{1}{k+1}\Bigr)
\]
(yang mungkin pada tiap langkahnya justru karena calonnya masih tersisa
tak hingga banyak) akan membangun barisan bagian yang konvergen ke $x$.
Bola $B(x, r_x)$ menyelimuti $X$; jika berhingga banyak di antaranya
menyelimuti $X$, maka himpunan indeks $\N$ akan menjadi gabungan hingga
atas himpunan hingga: mustahil.

($\Rightarrow$) Dua langkah. \emph{Bilangan Lebesgue:} untuk selimut
\omterm{def:b2:metric:topology}{terbuka} $(U_i)$ atas $X$ yang \omterm{def:b2:metric:compact}{kompak}, ada $\rho > 0$ sehingga setiap bola
berjari-jari $\rho$ terletak di suatu $U_i$. Sebab jika tidak, untuk tiap
$n$ pilihlah $x_n$ dengan $B(x_n, \frac{1}{n+1})$ tidak berada di $U_i$
mana pun; lalu sarikan $x_{\varphi(n)} \to x \in U_{i_0} \supseteq B(x,
r)$; untuk $n$ yang besar, $B(x_{\varphi(n)}, \frac{1}{\varphi(n)+1})
\subseteq B(x, r) \subseteq U_{i_0}$: kontradiksi. \emph{Keterbatasan
total:} untuk setiap $\rho > 0$, berhingga banyak bola berjari-jari
$\rho$ menyelimuti $X$. Sebab jika tidak, pilihlah secara induktif
$x_{n+1}$ di luar $B(x_0, \rho) \cup \dots \cup B(x_n, \rho)$: barisan
itu berjarak antara $\geq \rho$, jadi ia tidak punya barisan bagian yang
Cauchy --- sehingga tidak punya yang konvergen: kontradiksi. Setelah
digabungkan: selimuti $X$ oleh berhingga banyak bola berjari-jari $\rho$
(bilangan Lebesgue tadi), yang masing-masing berada di suatu $U_i$:
itulah subselimut hingganya.
\end{proof}

\begin{example}[Sebuah jaring-$\varepsilon$, dicacah]\label{ex:b2:metric:epsnet}
Keterbatasan total (dari bukti
\cref{thm:b2:metric:borellebesgue}) sangat konkret pada
$\intcc{0}{1}$: untuk $\varepsilon > 0$, sebanyak $\lceil
\frac{1}{2\varepsilon}\rceil$ bola berpusat di $\varepsilon,
3\varepsilon, 5\varepsilon, \dots$ dan berjari-jari $\varepsilon$
menyelimutinya --- kira-kira $\frac1{2\varepsilon}$ bola, dan tak ada
selimut yang sanggup dengan kurang dari $\frac{1}{2\varepsilon}$ di
antaranya (sebab tiap bola menyelimuti panjang paling banyak
$2\varepsilon$). Di $\intcc01^2$ cacahannya menjadi kuadratnya, yakni
berorde $\varepsilon^{-2}$: bilangan penyelimutan tumbuh seperti
$\varepsilon^{-d}$ pada dimensi $d$ --- sebuah wajah kuantitatif
kekompakan, sekaligus alasan mengapa bola satuan yang berdimensi tak hingga
pada \cref{ch:b2:nvs} (yang sama sekali tak punya jaring-$\frac13$
hingga) tak mungkin \omterm{def:b2:metric:compact}{kompak}.
\end{example}

\begin{example}[Membaca kekompakan pada selimut]\label{ex:b2:metric:covers}
Interval setengah \omterm{def:b2:metric:topology}{terbuka} $\intoc{0}{1}$ terselimuti oleh himpunan
\omterm{def:b2:metric:topology}{terbuka} $U_n = \intoo{\frac1n}{2}$, $n \geq 1$; sembarang keluarga bagian
yang hingga punya indeks terbesar $N$ dan melewatkan
$\intoc{0}{\frac1N}$: jadi tak ada subselimut hingga, sehingga
$\intoc{0}{1}$ tidak \omterm{def:b2:metric:compact}{kompak} --- dan definisi barisannya melihat hal itu
lewat $x_n = \frac1n$, yang limitnya $0$ meloloskan diri. Sebaliknya,
menambahkan satu titik $0$ itu saja memperbaiki kedua diagnosisnya sekaligus:
pada $\intcc{0}{1}$ setiap selimut semacam itu wajib memuat sebuah
himpunan yang memuat $0$, dan himpunan itu menelan satu ruas awal
seluruhnya, lalu berhingga banyak himpunan merampungkan sisanya. Kedua
bahasa \cref{thm:b2:metric:borellebesgue} selalu gagal atau berhasil
bersama-sama --- selimut mendeteksi pelarian persis di tempat barisan mendeteksinya.
\end{example}

\begin{remark}[Pandangan ke depan di dalam jilid ini]\label{rem:b2:metric:perspectives}
Bab ini adalah dinding pemikul jilid ini; perhatikan di mana
tiap pilarnya memikul beban. \emph{Kelengkapan}: kriteria
Cauchy menjadi uji kekonvergenan deret pada ruang Banach
(\cref{ch:b2:series}), kekonvergenan seragam pada
\cref{ch:b2:funcseq} tak lain kekonvergenan di dalam
$\bigl(C, d_\infty\bigr)$ yang \omterm{def:b2:metric:complete}{lengkap}, dan Cauchy--Lipschitz
(\cref{ch:b2:diffeq}) tak lain teorema titik tetap Banach yang
berbusana persamaan integral. \emph{Kekompakan}: ia membuktikan
kesetaraan norma (\cref{ch:b2:nvs}), tercapainya
ekstremum bagi optimasi pada \cref{ch:b2:diffcalc}, dan
keberadaan hampiran terbaik (pada soal akhir pekan
\cref{ch:b2:nvs}). \emph{Keterhubungan}: ia menggloballkan pernyataan
lokal --- ketunggalan penyelesaian persamaan diferensial, teorema
nilai antara pada kurva (\cref{ch:b2:curves}), serta
kedua komponen $GL_n(\R)$ yang kelak dipisahkan teori orientasi
(\cref{ch:b2:multint}).
\end{remark}

\begin{remark}[Jebakan yang sering muncul]\label{rem:b2:metric:pitfalls}
(i) ``Tertutup dan terbatas berarti \omterm{def:b2:metric:compact}{kompak}'' adalah teorema tentang
$\R^n$, bukan tentang \omterm{def:b2:metric:def}{ruang metrik}: himpunan tak hingga dengan
metrik diskret bersifat tertutup dan terbatas terhadap dirinya sendiri
namun tidak \omterm{def:b2:metric:compact}{kompak} (\cref{exo:b2:metric:4}), dan bola satuan tertutup
$C(\intcc01)$ pun gagal (\cref{ch:b2:nvs}). (ii) Kelengkapan adalah
sifat \emph{jaraknya}, bukan sifat topologinya: $\R$
dengan $d(x,y) = \abs{\arctan x - \arctan y}$ punya barisan konvergen
yang biasa tetapi tidak \omterm{def:b2:metric:complete}{lengkap}
(\cref{exo:b2:metric:1}). (iii) Bijeksi \omterm{def:b2:metric:continuity}{kontinu} belum tentu
homeomorfisma --- lihat parametrisasi lingkaran pada
\cref{exo:b2:metric:7}; kekompakan sumbernya memperbaikinya.
(iv) Teorema Banach memerlukan $k < 1$ secara \emph{seragam}: syarat
$d(f(x), f(y)) < d(x,y)$ semata tidak menjamin apa pun pada ruang
yang tidak \omterm{def:b2:metric:compact}{kompak} (\cref{exo:b2:metric:5}). (v) \omterm{def:b2:metric:connected}{Terhubung} tidak
berarti \omterm{def:b2:metric:connected}{terhubung} lintasan secara umum --- tetapi bagi himpunan
bagian \omterm{def:b2:metric:topology}{terbuka} ruang bernorma yang ditemui buku ini, keduanya sepadan
(\cref{ch:b2:nvs}).
\end{remark}

\begin{remark}[Di mana bab ini dipakai]
Di mana-mana pada paruh analisisnya. Kelengkapan
$C(\intcc{a}{b})$ menggerakkan teorema kekonvergenan pada
\cref{ch:b2:funcseq} dan teori Cauchy--Lipschitz pada
\cref{ch:b2:diffeq} (soal akhir pekan bab ini sudah membuktikan
teorema Picard--Lindelöf yang lokal); kekompakan memberi
kesetaraan norma dalam dimensi hingga (\cref{ch:b2:nvs}) dan
keberadaan ekstremum pada \cref{ch:b2:diffcalc}; sedangkan keterhubungan
melandasi hujah nilai antara pada \cref{ch:b2:realfun}
dan penggloballan ketunggalan bagi persamaan diferensial. Pada
jilid Tahun ke-3, kekompakan pada ruang fungsi (teorema
Arzelà--Ascoli) dan teorema kategori Baire
(\cref{exo:b2:metric:12} di sini) menjadi perkakas sehari-hari.
\end{remark}

\begin{omfigure}
\begin{tikzpicture}[xscale=6.4, yscale=0.8]
  \draw[omDef, line width=3pt] (0,3) -- (1,3);
  \draw[omDef, line width=3pt] (0,2) -- (0.3333,2);
  \draw[omDef, line width=3pt] (0.6667,2) -- (1,2);
  \foreach \a in {0, 0.6667} {
    \draw[omDef, line width=3pt] (\a,1) -- (\a + 0.1111,1);
    \draw[omDef, line width=3pt]
      (\a + 0.2222,1) -- (\a + 0.3333,1);
  }
  \foreach \a in {0, 0.2222, 0.6667, 0.8889} {
    \draw[omThm, line width=3pt] (\a,0) -- (\a + 0.037,0);
    \draw[omThm, line width=3pt]
      (\a + 0.0741,0) -- (\a + 0.1111,0);
  }
  \node[left, font=\small] at (0,3) {$C_0$};
  \node[left, font=\small] at (0,2) {$C_1$};
  \node[left, font=\small] at (0,1) {$C_2$};
  \node[left, font=\small] at (0,0) {$C_3$};
\end{tikzpicture}

{\small Tahap-tahap awal himpunan Cantor
(\cref{exo:b2:metric:8}): tiap tingkatnya membuang sepertiga tengah
yang \omterm{def:b2:metric:topology}{terbuka} dari setiap ruasnya. Irisan $C = \bigcap_n C_n$ bersifat
\omterm{def:b2:metric:compact}{kompak}, berinterior kosong dan berpanjang nol, namun \omterm{def:b2:structures:countable}{ekuipoten}
dengan $\R$ --- dan ia kembali sebagai \emph{titik tetap} sebuah
kontraksi atas himpunan pada soal akhir pekan bab ini (pertanyaan
22).}
\end{omfigure}

\section{Keterhubungan}

\begin{definition}\label{def:b2:metric:connected}
$X$ disebut \emph{terhubung}\index{ruang terhubung} apabila ia tidak
menerima pemilahan menjadi dua himpunan bagian \omterm{def:b2:metric:topology}{terbuka} yang tak kosong
--- setara dengan itu, apabila satu-satunya himpunan bagiannya yang
sekaligus \omterm{def:b2:metric:topology}{terbuka} dan tertutup adalah $\emptyset$ dan $X$. Adapun $X$
disebut \emph{terhubung lintasan} apabila sembarang dua titiknya
dihubungkan oleh pemetaan \omterm{def:b2:metric:continuity}{kontinu} $\gamma \colon \intcc{0}{1} \to X$.
\end{definition}

\begin{theorem}\label{thm:b2:metric:connectedness}
\begin{enumerate}
  \item Himpunan bagian $\R$ yang \omterm{def:b2:metric:connected}{terhubung} persis semua intervalnya.
  \item Peta \omterm{def:b2:metric:continuity}{kontinu} sebuah \omterm{def:b2:metric:connected}{ruang terhubung} bersifat \omterm{def:b2:metric:connected}{terhubung} ---
        dan dari sinilah teorema nilai antara yang umum: fungsi real
        \omterm{def:b2:metric:continuity}{kontinu} pada \omterm{def:b2:metric:connected}{ruang terhubung} mengambil setiap nilai di antara
        dua nilainya.
  \item \omterm{def:b2:metric:connected}{Terhubung} lintasan $\Rightarrow$ \omterm{def:b2:metric:connected}{terhubung}. (Konversnya gagal
        secara umum; ia berlaku bagi himpunan bagian \omterm{def:b2:metric:topology}{terbuka} ruang
        bernorma, \cref{ch:b2:nvs}.)
\end{enumerate}
\end{theorem}

\begin{proof}
(1) Himpunan $A$ yang bukan interval melewatkan suatu $z$ di antara dua
titiknya: maka $A = (A \cap \intoo{-\infty}{z}) \cup (A \cap
\intoo{z}{+\infty})$ memecahnya menjadi dua potongan tak kosong yang
\omterm{def:b2:metric:topology}{terbuka} (di $A$). Sebaliknya, misalkan $I$ sebuah interval dan $I = U
\cup V$ pemilahan menjadi himpunan tak kosong yang \omterm{def:b2:metric:topology}{terbuka} secara relatif;
pilih $a \in U$, $b \in V$, katakanlah $a < b$, lalu tetapkan
$s = \sup\,(U \cap \intcc{a}{b})$, sebuah titik $\intcc{a}{b}
\subseteq I$. Jika $s \in U$: maka $s \neq b$, dan keterbukaan relatif
$U$ menaruh satu interval penuh di sekitar $s$ (setelah diiriskan dengan
$I$) di dalam $U$ --- sehingga ada titik $U \cap \intcc{a}{b}$ yang
melampaui $s$, dan itu bertentangan dengan supremumnya. Jika $s \in V$:
keterbukaan relatif $V$ menaruh interval $\intoo{s - r}{s + r} \cap I$ di
dalam $V$; tetapi supremumnya melekat pada $U \cap \intcc{a}{b}$, yang
mestinya bertemu interval itu --- kontradiksi dengan $U \cap V =
\emptyset$. (Inilah hujah buka-tutup Tahun ke-1 untuk $\R$, dijalankan di dalam $I$.)

(2) Jika $f(X) = U' \cup V'$ terpecah menjadi himpunan tak kosong yang
\omterm{def:b2:metric:topology}{terbuka} secara relatif, maka $X = f^{-1}(U') \cup f^{-1}(V')$ memecah $X$
(\cref{thm:b2:metric:globalcontinuity}). Untuk teorema nilai antara:
$f(X) \subseteq \R$ \omterm{def:b2:metric:connected}{terhubung}, jadi ia interval menurut (1).

(3) Andaikan $X = U \cup V$ dengan keduanya \omterm{def:b2:metric:topology}{terbuka} dan tak kosong, lalu
hubungkan $a \in U$ ke $b \in V$ oleh lintasan $\gamma$: maka
$\gamma^{-1}(U), \gamma^{-1}(V)$ memecah $\intcc{0}{1}$, dan itu
bertentangan dengan (1).
\end{proof}

\begin{example}\label{ex:b2:metric:glnr}
$GL_n(\R)$ tidak \omterm{def:b2:metric:connected}{terhubung}: $\det$ bersifat \omterm{def:b2:metric:continuity}{kontinu} (sebuah polinomial
dalam unsurnya) ke $\R^*$, yang tidak \omterm{def:b2:metric:connected}{terhubung}; prapeta
$\R_+^*$ dan $\R_-^*$ memecah $GL_n(\R)$. (Tiap potongannya sebenarnya
\omterm{def:b2:metric:connected}{terhubung} lintasan --- latihan yang menyenangkan di luar kebutuhan
kita.) Sebaliknya $GL_n(\C)$ \emph{memang} \omterm{def:b2:metric:connected}{terhubung} lintasan:
\cref{exo:b2:metric:10}.
\end{example}

\begin{example}[Titik tetap dari keterhubungan semata]\label{ex:b2:metric:ivtfixed}
Setiap $f \colon \intcc01 \to \intcc01$ yang \omterm{def:b2:metric:continuity}{kontinu} punya titik
tetap --- tanpa hipotesis kontraksi, tanpa iterasi. Tinjaulah
$g(x) = f(x) - x$, yang \omterm{def:b2:metric:continuity}{kontinu} pada $\intcc01$ yang \omterm{def:b2:metric:connected}{terhubung}:
\[
g(0) = f(0) \geq 0,
\qquad
g(1) = f(1) - 1 \leq 0 ,
\]
lalu teorema nilai antara
(\cref{thm:b2:metric:connectedness}~(2)) menyerahkan sebuah nol bagi
$g$, yakni sebuah titik tetap bagi $f$. Bandingkan dengan Banach
(\cref{thm:b2:metric:banach}): di sini keberadaannya bersifat topologis
dan cuma-cuma, tetapi ketunggalan dan algoritmenya hilang --- sebab $f =
\mathrm{id}$ punya titik tetap di mana-mana, dan iterasi $f$ yang
bukan kontraksi bisa berputar selamanya. Kedua teorema titik tetap pada
bab ini menjawab pertanyaan yang berbeda dengan mata uang yang berbeda.
\end{example}

\begin{example}[$\R$ dan $\R^2$ tidak homeomorfik]\label{ex:b2:metric:rnotr2}
Keterhubungan adalah sidik jari topologis. Andaikan $h \colon
\R^2 \to \R$ sebuah homeomorfisma (bijeksi \omterm{def:b2:metric:continuity}{kontinu} yang inversnya juga
\omterm{def:b2:metric:continuity}{kontinu}). Buang satu titik $a \in \R^2$: maka
pembatasan $h \colon \R^2\setminus\{a\} \to
\R\setminus\{h(a)\}$ tetap homeomorfisma. Namun $\R^2$ dikurangi
satu titik itu bersifat \omterm{def:b2:metric:connected}{terhubung} lintasan --- hubungkan sembarang dua
titiknya oleh ruas garis, dengan memutar lewat ruas kedua melalui titik
bantu bila $a$ menghalangi ruas langsungnya --- jadi ia \omterm{def:b2:metric:connected}{terhubung}
(\cref{thm:b2:metric:connectedness}~(3)); sedangkan $\R$ dikurangi
satu titik terpecah menjadi dua setengah-garis \omterm{def:b2:metric:topology}{terbuka} yang tak kosong:
tidak \omterm{def:b2:metric:connected}{terhubung}. Padahal keterhubungan terpelihara oleh pemetaan \omterm{def:b2:metric:continuity}{kontinu}:
kontradiksi. Bidang dan garis sungguh berbeda \emph{sebagai
ruang topologis} --- kenyataan yang terlalu kasar untuk dilihat
kardinalitas semata (bijeksi bergaya
\cref{exo:b2:structures:3} memang ada!).
\end{example}

\section{Latihan}

\begin{exercise}[$\star$]\label{exo:b2:metric:1}
Pada $\R$, periksalah bahwa $\delta(x, y) = \min(1, \abs{x - y})$ dan
$d(x,y) = \abs{\arctan x - \arctan y}$ merupakan jarak. Barisan mana yang
konvergen untuk masing-masingnya? Apakah $(\R, d)$ \omterm{def:b2:metric:complete}{lengkap}?
\end{exercise}

\begin{exercise}[$\star$]\label{exo:b2:metric:2}
Pada sebuah \omterm{def:b2:metric:def}{ruang metrik}, buktikan bahwa barisan konvergen bersifat
Cauchy dan terbatas, dan bahwa barisan Cauchy yang punya barisan bagian
konvergen pastilah konvergen. Turunkan sekali lagi bahwa \omterm{def:b2:metric:compact}{ruang metrik
kompak} bersifat \omterm{def:b2:metric:complete}{lengkap}.
\end{exercise}

\begin{exercise}[$\star$]\label{exo:b2:metric:3}
Di $\bigl(C(\intcc{0}{1}), d_\infty\bigr)$, hitunglah jarak
antara $f(x) = x$ dan $g(x) = x^2$; gambarkan bola tertutup
$\overline B(0, 1)$; lalu buktikan bahwa himpunan $\{f : f(0) = 0\}$
tertutup sedangkan $\{f : f(0) > 0\}$ \omterm{def:b2:metric:topology}{terbuka}.
\end{exercise}

\begin{exercise}[$\star\star$]\label{exo:b2:metric:4}
Buktikan bahwa \omterm{def:b2:metric:def}{ruang metrik} diskret $X$ (himpunan apa pun) bersifat
\omterm{def:b2:metric:complete}{lengkap}, dan bahwa ia \omterm{def:b2:metric:compact}{kompak} bila dan hanya bila $X$ hingga. Himpunan
bagian mana yang \omterm{def:b2:metric:connected}{terhubung}?
\end{exercise}

\begin{exercise}[$\star\star$]\label{exo:b2:metric:5}
Misalkan $X$ \omterm{def:b2:metric:compact}{kompak} dan $f \colon X \to X$ dengan
\[
d\bigl(f(x), f(y)\bigr) < d(x, y) \quad \text{untuk setiap } x \neq y .
\]
Buktikan bahwa $f$ punya tepat satu titik tetap \emph{(minimumkan $x
\mapsto d(x, f(x))$)}, lalu berikan contoh pada $X = \intco{1}{+\infty}$
(yang tidak \omterm{def:b2:metric:compact}{kompak}) tanpa titik tetap.
\end{exercise}

\begin{exercise}[$\star\star$]\label{exo:b2:metric:6}
(Himpunan \omterm{def:b2:metric:compact}{kompak} bersarang) Misalkan $(K_n)$ barisan turun berisi
himpunan bagian \omterm{def:b2:metric:compact}{kompak} yang tak kosong pada sebuah \omterm{def:b2:metric:def}{ruang metrik}.
Buktikan bahwa $\bigcap_n K_n \neq \emptyset$ \emph{(pilih $x_n \in K_n$
lalu sarikan)}. Tunjukkan lewat sebuah contoh bahwa himpunan
\emph{tertutup} tak kosong yang bersarang di $\R$ dapat beririsan kosong.
\end{exercise}

\begin{exercise}[$\star\star$]\label{exo:b2:metric:7}
Misalkan $K$ \omterm{def:b2:metric:compact}{kompak} dan $f \colon K \to Y$ \omterm{def:b2:metric:continuity}{kontinu} serta bijektif.
Buktikan bahwa $f^{-1}$ \omterm{def:b2:metric:continuity}{kontinu} \emph{(pakailah himpunan tertutup:
\cref{thm:b2:metric:globalcontinuity} dan
\cref{thm:b2:metric:compactprops})}. Berikan contoh penyangkal tanpa
kekompakan ($\gamma(t) = (\cos t, \sin t)$ pada
$\intco{0}{2\pi}$).
\end{exercise}

\begin{exercise}[$\star\star$]\label{exo:b2:metric:8}
\emph{Himpunan Cantor} $C$ diperoleh dari $\intcc{0}{1}$ dengan
membuang sepertiga tengah yang \omterm{def:b2:metric:topology}{terbuka} berulang kali. Buktikan bahwa $C$
\omterm{def:b2:metric:compact}{kompak}, berinterior kosong, dan tak hingga --- bahkan \omterm{def:b2:structures:countable}{ekuipoten} dengan
$\{0,1\}^{\N}$ \emph{(lewat uraian terner berangka $0,2$;
\cref{exo:b2:structures:3})}.
\end{exercise}

\begin{exercise}[$\star\star\star$]\label{exo:b2:metric:9}
Misalkan $X$ \omterm{def:b2:metric:def}{ruang metrik} \emph{\omterm{def:b2:metric:compact}{kompak}} dan $f \colon X \to X$ sebuah
isometri: $d(f(x), f(y)) = d(x, y)$. Buktikan bahwa $f$ surjektif.
\emph{Petunjuk: jika $a \notin f(X)$, maka $\varepsilon = d(a, f(X)) > 0$
(mengapa?); telaahlah orbit $a, f(a), f^2(a), \dots$ lalu tunjukkan
titiknya berjarak $\geq \varepsilon$ sepasang demi sepasang --- dan itu
bertentangan dengan kekompakan.}
\end{exercise}

\begin{exercise}[$\star\star\star$]\label{exo:b2:metric:10}
Buktikan bahwa $GL_n(\C)$ \omterm{def:b2:metric:connected}{terhubung} lintasan. \emph{Petunjuk: diberikan
$A, B$ yang punya invers, tinjaulah $p(z) = \det\bigl((1 - z)A + zB\bigr)$
untuk $z \in \C$: ia polinomial dalam $z$ yang tidak nol secara identik,
jadi ia punya berhingga banyak akar; hubungkan $0$ ke $1$ di $\C$ lewat
lintasan yang menghindari akar itu.}
\end{exercise}

\begin{exercise}[$\star\star$]\label{exo:b2:metric:11}
Untuk $\emptyset \neq A \subseteq X$, tetapkan $d(x, A) = \inf_{a \in A}
d(x, a)$. Buktikan bahwa $x \mapsto d(x, A)$ bersifat \omterm{def:b2:metric:continuity}{Lipschitz} dengan
konstanta $1$, bahwa $d(x, A) = 0$ bila dan hanya bila $x \in \overline
A$, dan bahwa untuk himpunan \emph{tertutup} $A, B$ yang tak kosong dan
saling lepas, fungsi
\[
\varphi(x) = \frac{d(x, A)}{d(x, A) + d(x, B)}
\]
terdefinisi dengan baik, \omterm{def:b2:metric:continuity}{kontinu}, bernilai $0$ persis pada $A$ dan
bernilai $1$ persis pada $B$ --- yakni sebuah ``sakelar'' \omterm{def:b2:metric:continuity}{kontinu} yang
memisahkan sembarang dua himpunan tertutup yang saling lepas.
\end{exercise}

\begin{exercise}[$\star\star\star$]\label{exo:b2:metric:12}
(Baire) Misalkan $X$ \omterm{def:b2:metric:complete}{ruang metrik lengkap} dan $(U_n)_{n\geq1}$
barisan himpunan bagian \omterm{def:b2:metric:topology}{terbuka} yang padat. Buktikan bahwa $\bigcap_n
U_n$ padat di $X$ \emph{(di dalam sembarang bola, bangunlah bola
tertutup bersarang $\overline B(x_n, r_n) \subseteq U_n$ dengan $r_n \to
0$ lalu pakai kelengkapannya)}. Turunkan bahwa $\R$ bukan gabungan
\omterm{def:b2:structures:countable}{terbilang} atas himpunan tertutup yang berinterior kosong, lalu peroleh
kembali --- sekali lagi --- bahwa $\R$ tidak \omterm{def:b2:structures:countable}{terbilang}.
\end{exercise}

\section{Soal: Iterasi Picard}

Kelengkapan ditambah kontraksi adalah mesin penyelesai: berilah ia sebuah
persamaan yang ditulis sebagai masalah titik tetap, dan ia mengembalikan
keberadaan, ketunggalan, algoritme, beserta batang galatnya. Soal akhir
pekan ini menjalankan mesin itu dengan tenaga penuh pada persamaan $y' =
f(t, y)$: kita membuktikan \emph{teorema Picard--Lindelöf} yang lokal
(yakni jantung taklinear teori Cauchy--Lipschitz pada
\cref{ch:b2:diffeq}), menyaksikan tiap hipotesisnya membayar sewa lewat
contoh penyangkal, lalu memungut panen yang murni metrik ---
kebergantungan \omterm{def:b2:metric:continuity}{kontinu} pada datanya, persamaan Kepler, dan
keserupaan-diri himpunan Cantor.

\begin{omfigure}
\begin{tikzpicture}
\begin{axis}[omaxis, width=8.6cm, height=5.4cm,
    xmin=-0.05, xmax=1.45, ymin=0, ymax=4.2,
    xtick={0,1}, ytick={1,2,3}]
  \addplot[omDef!45, thick, domain=0:1.4, samples=2] {1};
  \addplot[omDef!65, thick, domain=0:1.4, samples=2] {1 + x};
  \addplot[omDef, thick, domain=0:1.4, samples=60]
    {1 + x + x^2/2};
  \addplot[omThm, very thick, domain=0:1.4, samples=120]
    {exp(x)};
  \node[font=\small, omDef] at (axis cs:1.28,0.8) {$y_0$};
  \node[font=\small, omDef] at (axis cs:1.28,2.15) {$y_1$};
  \node[font=\small, omDef] at (axis cs:1.28,2.95) {$y_2$};
  \node[font=\small, omThm] at (axis cs:1.17,3.85)
    {$\eu^{\,t}$};
\end{axis}
\end{tikzpicture}

{\small Iterat Picard bagi $y' = y$, $y(0) = 1$: tiap lintasan
lewat $T(y)(t) = 1 + \int_0^t y$ menambahkan satu suku Taylor, dan
kontraksinya memeras seluruh barisan itu secara seragam ke
$\eu^{\,t}$.}
\end{omfigure}

\begin{problem}[{Soal akhir pekan --- teorema
Picard--Lindelöf}]
\label{pb:b2:metric:1}
Di sepanjang soal ini, $t_0 \in \R$, $y_0 \in \R$, $a, b > 0$, dan $f$
adalah fungsi \omterm{def:b2:metric:continuity}{kontinu} pada persegi panjang $R = \intcc{t_0 - a}{t_0 +
a} \times \intcc{y_0 - b}{y_0 + b}$, yang terbatas oleh $M =
\sup_R\,\abs f$ dan \emph{Lipschitz-$L$ pada peubah keduanya}: yakni
$\abs{f(t, y) - f(t, z)} \leq L\abs{y - z}$ setiap kali kedua titiknya
berada di $R$. Tetapkan
\[
h = \min\Bigl(a, \frac bM\Bigr) \quad (\text{dengan } h = a
\text{ bila } M = 0), \qquad I = \intcc{t_0 - h}{t_0 + h}.
\]

\medskip
\noindent\textbf{Bagian I --- Panggung yang \omterm{def:b2:metric:complete}{lengkap}.}
\begin{enumerate}
  \item Buktikan kedua pernyataan yang dikutip pada
        \cref{def:b2:metric:complete}: himpunan bagian tertutup dari
        \omterm{def:b2:metric:complete}{ruang metrik lengkap} bersifat \omterm{def:b2:metric:complete}{lengkap}, dan himpunan bagian
        \omterm{def:b2:metric:complete}{lengkap} dari \omterm{def:b2:metric:def}{ruang metrik} apa pun bersifat tertutup. Turunkan
        bahwa setiap himpunan bagian tertutup $\bigl(C(I),
        d_\infty\bigr)$ merupakan \omterm{def:b2:metric:def}{ruang metrik} yang \omterm{def:b2:metric:complete}{lengkap}.
  \item Tunjukkan bahwa pemetaan $C(I) \to C(I)$, $y \mapsto \bigl(t
        \mapsto \int_{t_0}^{t}y(s)\,\dd s\bigr)$, bersifat
        Lipschitz-$h$ terhadap $d_\infty$.
  \item (Titik tetap bergerak lebih sedikit daripada pemetaannya)
        Misalkan $g \colon X \to X$ kontraksi-$k$ sebuah \omterm{def:b2:metric:def}{ruang metrik}
        dengan titik tetap $\ell_g$, dan $\widetilde g \colon X \to X$
        pemetaan \emph{apa pun} yang punya titik tetap
        $\ell_{\widetilde g}$. Buktikan
        \[
        d(\ell_g, \ell_{\widetilde g}) \leq
        \frac{d\bigl(g(\ell_{\widetilde g}),
        \widetilde g(\ell_{\widetilde g})\bigr)}{1 - k}
        \leq \frac{\sup_{x \in X} d\bigl(g(x), \widetilde
        g(x)\bigr)}{1 - k}.
        \]
  \item (Siasat iterat) Misalkan $X$ \omterm{def:b2:metric:complete}{lengkap} dan tak kosong serta $g
        \colon X \to X$ sebuah pemetaan --- yang tidak diandaikan
        \omterm{def:b2:metric:continuity}{kontinu} --- sedemikian sehingga suatu iterat $g^m$ merupakan
        kontraksi-$k$. Buktikan bahwa $g$ punya tepat satu titik tetap
        $\ell$ dan bahwa \emph{setiap} orbit $x_{n+1} = g(x_n)$
        konvergen ke $\ell$. \emph{(Titik tetap $g$ adalah titik tetap
        $g^m$; sebaliknya $g(\ell)$ merupakan titik tetap $g^m$;
        pecahlah orbitnya menurut sisa modulo $m$.)}
\end{enumerate}

\medskip
\noindent\textbf{Bagian II --- Teorema Picard--Lindelöf.}
\begin{enumerate}[resume]
  \item Tunjukkan bahwa fungsi $y \colon I \to \intcc{y_0 - b}{y_0
        + b}$ bersifat $C^1$ dengan $y(t_0) = y_0$ dan $y' = f(t, y)$
        pada $I$ bila dan hanya bila ia \omterm{def:b2:metric:continuity}{kontinu} dan memenuhi
        persamaan integral
        \[
        y(t) = y_0 + \int_{t_0}^{t} f\bigl(s, y(s)\bigr)\dd s
        \qquad (t \in I).
        \]
  \item Misalkan $X_h = \{y \in C(I) : \abs{y(t) - y_0} \leq b
        \text{ pada } I\}$ dan misalkan $T$ ditetapkan oleh $T(y)(t) =
        y_0 + \int_{t_0}^{t}f(s, y(s))\dd s$. Tunjukkan bahwa $X_h$
        merupakan himpunan bagian tertutup $C(I)$ yang tak kosong, jadi
        \omterm{def:b2:metric:complete}{lengkap}, dan bahwa $T$ memetakan $X_h$ ke dalam $X_h$ ---
        di sinilah $h \leq b/M$ bekerja.
  \item Tunjukkan bahwa $d_\infty\bigl(T(y), T(z)\bigr) \leq
        Lh\,d_\infty(y, z)$ pada $X_h$: jika $Lh < 1$, teorema Banach
        sudah merampungkannya. Syarat kekecilan itu kita singkirkan berikutnya.
  \item Buktikan secara induktif pada $n$:
        \[
        \abs{T^n(y)(t) - T^n(z)(t)} \leq
        \frac{\bigl(L\abs{t - t_0}\bigr)^n}{n!}\,
        d_\infty(y, z) \qquad (y, z \in X_h,\ t \in I),
        \]
        sehingga suatu iterat $T$ merupakan kontraksi. Rampungkan dengan
        pertanyaan 4 (\emph{teorema Picard--Lindelöf}): masalah
        Cauchy $y' = f(t,y)$, $y(t_0) = y_0$ punya tepat
        satu penyelesaian pada $I = \intcc{t_0 - h}{t_0 + h}$ yang
        bernilai di $\intcc{y_0 - b}{y_0 + b}$.
  \item Tunjukkan bahwa pembatasan ``yang bernilai di $\intcc{y_0 -
        b}{y_0 + b}$'' bersifat otomatis: setiap penyelesaian masalah
        Cauchy yang terdefinisi pada $I$ dan grafiknya berawal di $R$
        akan tetap berada di $\intcc{y_0 - b}{y_0 + b}$ \emph{(tinjaulah
        saat keluar yang pertama lalu batasi $\abs{y(t) - y_0}$ oleh
        $M\abs{t - t_0}$)}. Karenanya ketunggalannya berlaku di antara
        semua penyelesaian pada $I$.
  \item Jalankan mesinnya pada $y' = y$, $y(0) = 1$, berawal dari
        fungsi konstan $y^{(0)} \equiv 1$: hitunglah iterat
        Picard $y^{(n)}$, kenali bentuknya, lalu gambarkan
        kekonvergenannya.
\end{enumerate}

\medskip
\noindent\textbf{Bagian III --- Tiap hipotesis membayar sewanya.}
\begin{enumerate}[resume]
  \item (\omterm{def:b2:metric:continuity}{Lipschitz} gagal, ketunggalan gagal) Untuk $y' =
        2\sqrt{\abs y}$, $y(0) = 0$: periksalah bahwa $y \equiv 0$
        dan, untuk setiap $c \geq 0$, fungsi $y_c(t) = 0$ untuk
        $t \leq c$ serta $y_c(t) = (t - c)^2$ untuk $t > c$, semuanya
        merupakan penyelesaian $C^1$ pada $\R$. Di mana persisnya $y
        \mapsto 2\sqrt{\abs y}$ gagal menjadi \omterm{def:b2:metric:continuity}{Lipschitz}?
  \item (Kelokalannya nyata) Untuk $y' = y^2$, $y(0) = 1$: selesaikan
        secara gamblang, berikan interval keberadaan yang maksimal, lalu
        hitunglah $h$ terbaik yang dapat disahkan teoremanya atas semua
        pilihan persegi panjangnya ($a$ besar, $b$ bebas): tunjukkan
        $h_{\max} = \sup_{b>0} \frac{b}{(1+b)^2} = \frac14$,
        padahal penyelesaian sejatinya hidup pada $\intoo{-\infty}{1}$.
  \item (Kelengkapan bukan hiasan) Pada $X = \Q \cap
        \intcc{1}{2}$ dengan jarak biasa, misalkan $g(x) =
        \frac x2 + \frac1x$. Tunjukkan $g(X) \subseteq X$, bahwa $g$
        merupakan kontraksi-$\frac12$ \emph{(lewat ketaksamaan nilai
        rata-rata)}, dan bahwa $g$ tidak punya titik tetap di $X$.
        Hipotesis teorema Banach yang manakah yang gagal, dan apa
        titik tetapnya di dalam pelengkapannya?
  \item (Batang galat) Untuk kontraksi-$k$ bernama $g$ pada ruang yang
        \omterm{def:b2:metric:complete}{lengkap}, buktikan taksiran \emph{aposteriori} $d(x_n,
        \ell) \leq \frac{k}{1-k}\,d(x_n, x_{n-1})$. Untuk pemetaan
        Heron $g(x) = \frac x2 + \frac 1x$ pada $\intcc{1}{2}$
        (dengan titik tetap $\sqrt2$), berawal dari $x_0 = \frac32$:
        berapa langkah yang dituntut batas \emph{apriori}
        $\frac{k^n}{1-k}d(x_1, x_0)$ untuk ketelitian
        $10^{-6}$, dan berapa langkah yang sesungguhnya cukup?
        (Hitunglah $x_1, x_2, x_3$ beserta galatnya; batas
        kontraksinya jujur tetapi pesimistis --- sebab
        Heron konvergen secara kuadratik.)
\end{enumerate}

\medskip
\noindent\textbf{Bagian IV --- Kebergantungan \omterm{def:b2:metric:continuity}{kontinu}.} Pada bagian
ini $Lh < 1$, jadi $T$ sendiri merupakan kontraksi pada $X_h$ (pertanyaan
7); tulis $y[\,y_0\,]$ untuk penyelesaian yang bernilai awal
$y_0$.
\begin{enumerate}[resume]
  \item (Kebergantungan pada nilai awal) Misalkan $z_0$ nilai awal
        lain dengan $\abs{z_0 - y_0}$ cukup kecil sehingga
        kedua masalahnya muat di persegi panjangnya. Dengan memakai
        pertanyaan 3, buktikan
        \[
        d_\infty\bigl(y[y_0], y[z_0]\bigr) \leq
        \frac{\abs{y_0 - z_0}}{1 - Lh}.
        \]
  \item (Interval panjang lewat perangkaian) Andaikan penyelesaiannya
        ada pada sebuah ruas panjang yang dipotong menjadi $m$ potongan
        berurutan, dan pada masing-masingnya batas tadi berlaku dengan
        $Lh \leq \frac12$. Tunjukkan bahwa simpangannya tumbuh paling
        banyak dengan faktor $2$ tiap potongan, sehingga $d_\infty \leq
        2^m\abs{y_0 - z_0}$ secara keseluruhan --- sebuah batas yang
        eksponensial terhadap panjangnya, yakni bayangan diskret
        $\eu^{L\abs{t - t_0}}$ pada lema Gronwall
        (\cref{ch:b2:diffeq}).
  \item (Kebergantungan pada medannya) Misalkan $g$ medan lain pada
        $R$, yang juga Lipschitz-$L$ pada $y$, dengan $\sup_R \abs{f - g}
        \leq \varepsilon$. Buktikan bahwa penyelesaian yang
        bersesuaian memenuhi $d_\infty \leq
        \frac{\varepsilon h}{1 - Lh}$: jadi galat permodelannya
        merambat secara linear.
  \item (Parameter) Jika keluarga medan $f_\lambda$ bersifat
        Lipschitz-$L$ pada $y$ secara seragam dan $\sup_R\abs{f_\lambda
        - f_\mu} \leq C\abs{\lambda - \mu}$, turunkan bahwa
        $\lambda \mapsto y_\lambda$ bersifat \omterm{def:b2:metric:continuity}{Lipschitz} dari ruang
        parameternya ke $\bigl(C(I), d_\infty\bigr)$.
  \item (Sistem tidak berbiaya apa pun) Jelaskan mengapa Bagian I, II
        dan IV berlaku kata demi kata untuk $y$ yang bernilai di $\R^n$
        (dengan jarak supremum yang \omterm{def:b2:structures:generated}{dibangun} atas salah satu jarak pada
        \cref{ex:b2:metric:examples}), lalu hitunglah semua iterat
        Picard untuk sistem $y' = Ay$, $y(0) = (c_1,
        c_2)$, $A = \begin{psmallmatrix} 0 & 1\\ 0 &
        0\end{psmallmatrix}$: tunjukkan bahwa iterasinya menjadi
        mandek pada penyelesaian eksaknya setelah satu langkah.
\end{enumerate}

\medskip
\noindent\textbf{Bagian V --- Panen metrik dan rangkuman.}
\begin{enumerate}[resume]
  \item (Usikan identitas) Misalkan $X$ panggung \omterm{def:b2:metric:complete}{lengkap} yang
        bergaya ruang bernorma: ambil $X = C(I)$ atau $\R^n$. Jika
        $\eta \colon X \to X$ bersifat Lipschitz-$k$ dengan $k < 1$,
        buktikan bahwa $x \mapsto x + \eta(x)$ merupakan bijeksi $X$
        yang inversnya Lipschitz-$\frac1{1-k}$ \emph{(untuk tiap
        $y$, terapkan Banach pada $x \mapsto y - \eta(x)$)}. Inilah
        jantung metrik teorema fungsi invers
        (\cref{ch:b2:diffcalc}).
  \item (Persamaan Kepler) Untuk $0 \leq e < 1$ dan $m \in \R$,
        buktikan bahwa $x = m + e\sin x$ punya tepat satu penyelesaian,
        bahwa iterasi $x_{n+1} = m + e\sin x_n$ konvergen
        ke sana dari awal mana pun, lalu taksirlah: untuk $e = \frac12$,
        $m = 1$, berapa iterasi yang menjamin galat $\leq
        10^{-3}$ menurut batas apriorinya? (Penyelesaiannya $x
        \approx 1.4987$.)
  \item (Himpunan Cantor sebagai titik tetap) Misalkan $S_1(x) = \frac
        x3$ dan $S_2(x) = \frac x3 + \frac23$ pada $\R$, dan misalkan
        $C$ himpunan Cantor pada \cref{exo:b2:metric:8}. Buktikan
        $C = S_1(C) \cup S_2(C)$, lalu jelaskan dalam satu kalimat
        mengapa tak ada himpunan \omterm{def:b2:metric:compact}{kompak} tak kosong \emph{lain} yang
        memenuhi persamaan itu (pemetaan $A \mapsto S_1(A) \cup S_2(A)$
        merupakan kontraksi terhadap sebuah jarak antara himpunan
        \omterm{def:b2:metric:compact}{kompak} --- yaitu jarak Hausdorff, yang dijujurkan pada jilid Tahun ke-3).
  \item (Keterhubungan menggloballkan ketunggalan) Misalkan $f$ bersifat
        \omterm{def:b2:metric:continuity}{Lipschitz} secara lokal pada $y$ di sebuah himpunan \omterm{def:b2:metric:topology}{terbuka}, dan
        misalkan $y, z$ dua penyelesaian $y' = f(t, y)$ pada interval
        bersama $J$ dengan $y(t_0) = z(t_0)$. Buktikan $y = z$ pada $J$:
        tunjukkan bahwa $\{t \in J : y(t) = z(t)\}$ tak kosong, tertutup
        di $J$, dan \omterm{def:b2:metric:topology}{terbuka} di $J$ (lewat ketunggalan lokalnya), lalu
        pakailah keterhubungan interval
        (\cref{thm:b2:metric:connectedness}).
  \item (Tanpa kekecilan bagi persamaan linear) Untuk $y' =
        \alpha(t)y + \beta(t)$ dengan $\alpha, \beta$ \omterm{def:b2:metric:continuity}{kontinu}
        pada ruas $\intcc{A}{B}$, sesuaikan pertanyaan 8 untuk menunjukkan
        bahwa batas faktorialnya berlaku pada \emph{seluruh} ruas itu,
        sehingga keberadaan dan ketunggalannya bersifat global di sana ---
        yakni kasus skalar teorema Cauchy--Lipschitz pada
        \cref{ch:b2:diffeq}, tanpa pembatasan pada panjang $B - A$.
  \item (Rangkuman) Satu kalimat untuk masing-masing: apa yang
        disumbangkan kelengkapan; apa yang disumbangkan kontraksinya; apa
        yang dibeli siasat iterat dibandingkan Banach biasa; di mana
        keterhubungan masuk; dan contoh penyangkal Bagian
        III yang mana menjaga hipotesis yang mana. Sebutkan teorema
        puncaknya, lalu katakan apa yang menggantikan kontraksi ketika $f$
        sekadar \omterm{def:b2:metric:continuity}{kontinu} (teorema Peano, lewat kekompakan pada ruang
        fungsi --- yaitu Arzelà--Ascoli pada jilid Tahun ke-3).

\end{enumerate}
\end{problem}
